Визуализация течения жидкости \textemdash{} одна из классических задач компьютерной графики. Впервые она была решена еще в 80-х, но тем не менее реалистичное моделирование все еще остается весьма сложной задачей, совмещающей в себе непростую теоретическую часть и большой объем математических вычислений. Стоит отметить, что под симуляцией жидкости могут понимать как гидро, так и газодинамических процессы, например вода или дым. С практической точки зрения, симуляция жидкости может быть полезна в самых разных областях \textemdash{} от спецэффектов в кино и игровых движков до инженерных расчётов аэродинамики и анализах атмосферных явлений. Причем в прикладных задачах важны не столько строгость физических моделей, сколько сохранение ключевых качеств течения (несжимаемость, вихревая структура, гладкость поля скоростей) при приемлемой вычислительной стоимости.

Физическая основа моделирования обычно строится на уравнениях Навье–Стокса, которые формулируют законы сохранения массы и импульса. Полное решение этих уравнений требует больших вычислительных ресурсов, поэтому в задачах интерактивной визуализации применяются устойчивые приближённые схемы, дающие адекватную визуальную достоверность при существенно меньших затратах вычислений.

С точки зрения компьютерных систем, современные методы вычислений ориентированы на перенос интенсивных расчетов на GPU, а сами симуляции зачастую разрабатываются не под конкретные ОС, а сразу под несколько. Относительно современным решением последней проблемы стала поддержка браузерами методов рендеринга на GPU, что позволяет достичь кросс-платформенной совместимости по определению. Долгое время в вебе доминировал WebGL, основанный на классическом графическом конвейере. Однако появление WebGPU изменило ситуацию. Этот модуль предоставил более низкоуровневый доступ, унифицированные вычислительные шейдеры и явное управление ресурсами, что довольно близко по модели к Vulkan, Direct3D~12 и Metal. Это делает реализацию численных алгоритмов в браузере более эффективной и переносимой между платформами.

В данной курсовой работе рассматривается интерактивная визуализация двумерного течения жидкости с использованием WebGPU. В качестве платформы выбрано веб-приложение на TypeScript с использованием фреймворка Next.js: такой стек обеспечивает упрощенный процесс разработки самого приложения и позволяет эффективно задействовать возможности GPU через WebGPU для выполнения вычислительно сложных операций.

Цель работы \textemdash{} исследовать алгоритмы и разработать веб-приложение, демонстрирующее визуализацию течения жидкости в реальном времени на основе выбранного численного алгоритма. В рамках работы планируется описать физико-математические основы моделирования, обосновать выбор алгоритма и технологий, реализовать и постараться оптимизировать решение с точки зрения качества визуализации и производительности. Результатом станет интерактивное приложение, позволяющее пользователю наблюдать и взаимодействовать с моделью течения жидкости.
